% GENERATED, edit the French unit or the translation table
\mxunit{automatismes-ch02}{Lines in the plane}

\begin{fluencysheet}{Mon~2~Nov.~2026}{Direction vector of a line}
\mxfq{Warm-up · Quick}{GEO}{Give a vector collinear to $\vec u\vc{2}{-4}$ with simpler integer coordinates.}{$\vc{1}{-2}$ (divide by $2$).}{}
\mxfq{Core · Method}{GEO}{Give a direction vector of $(AB)$ for $A\pt{1;2}$ and $B\pt{4;8}$.}{$\vect{AB}\vc{3}{6}$, or $\vc{1}{2}$, which is collinear to it.}{}
\mxfq{Core + · MCQ}{GEO}{A direction vector of the line through $A\pt{0;1}$ and $B\pt{2;-3}$ is:\mxchoices{$\vc{1}{-2}$}{$\vc{1}{-4}$}{$\vc{-4}{-2}$}{$\vc{2}{4}$}}{A.\ $\vc{1}{-2}$}{Here $\vect{AB}\vc{2}{-4}$; test collinearity with the determinant.}
\mxfq{Review · Quick}{STA}{Marks: $8,\,12,\,12,\,15,\,18$. Calculate the mean, the median and the range.}{Mean $13$; median $12$; range $10$.}{}
\mxfq{Thread · Property}{DEM}{$ABC$ is right-angled at $A$ and $I$ is the midpoint of $[BC]$. What can be said about $IA$, $IB$ and $IC$? Why?}{$IA=IB=IC$: the midpoint of the hypotenuse is the centre of the circumcircle.}{}
\end{fluencysheet}

\begin{fluencysheet}{Tue~3~Nov.~2026}{Cartesian equation of a line}
\mxfq{Warm-up · Quick}{SUB}{Calculate $3x-2y+1$ for $x=1$ and $y=-2$.}{$3+4+1=8$.}{}
\mxfq{Core · Method}{GEO}{The line $d$ has equation $2x+3y-6=0$. Give a direction vector and a point of $d$.}{For $ax+by+c=0$, $\vc{-b}{a}$ is a direction vector: $\vc{-3}{2}$. The point $\pt{3;0}$ works: $6+0-6=0$.}{}
\mxfq{Core + · True/false}{GEO}{True or false? Justify. The lines $2x+3y-6=0$ and $4x+6y+1=0$ are parallel.}{True. Direction vectors $\vc{-3}{2}$ and $\vc{-6}{4}$ are collinear; the lines are distinct: $\pt{3;0}$ does not lie on the second line ($12+1\neq0$).}{}
\mxfq{Review · Table}{LEC}{Mode of transport of $30$ students.\par\smallskip\mxarraystretch{1.15}\begin{tabular}{|c|c|c|c|c|}\hline
Transport & Bus & Bike & Walking & Car \\ \hline
Frequency & 12 & 5 & 8 & 5 \\ \hline
\end{tabular}\par Give the proportion of students who come by bus and the angle of the ``bike'' sector in a pie chart.}{$\dfrac{12}{30}=\pc{40}$; $\dfrac{5}{30}\times360\degre=60\degre$.}{}
\mxfq{Thread · Quick}{REP}{On a number line, $A$ has coordinate $-2$ and $B$ has coordinate $3$. Give $AB$ and the coordinate of the midpoint of $[AB]$.}{$AB=3-(-2)=5$; midpoint: $\dfrac{-2+3}{2}=\num{0,5}$.}{}
\end{fluencysheet}

\begin{fluencysheet}{Thu~5~Nov.~2026}{Slope-intercept equation and gradient of a non-vertical line}
\mxfq{Warm-up · Quick}{EQ}{Make $y$ the subject of $2x+y-5=0$.}{$y=-2x+5$.}{}
\mxfq{Core · Method}{GEO}{Give the gradient and the y-intercept of the line $3x-2y+4=0$.}{$-2y=-3x-4$, $y=\dfrac32x+2$: gradient $\dfrac32$, y-intercept $2$.}{Pitfall: divide every term by $-2$ (signs).}
\mxfq{Core + · MCQ}{GEO}{The gradient of the line $(AB)$ with $A\pt{1;2}$ and $B\pt{5;0}$ is:\mxchoices{$2$}{$\dfrac12$}{$-2$}{$-\dfrac12$}}{D.\ $-\dfrac12$}{Gradient $=\dfrac{y_B-y_A}{x_B-x_A}=\dfrac{-2}{4}$.}
\mxfq{Review · Quick}{CAL}{Calculate $\dfrac34+\dfrac56$.}{$\dfrac{9}{12}+\dfrac{10}{12}=\dfrac{19}{12}$.}{}
\mxfq{Thread · Problem}{DIV}{Two gears have $12$ and $18$ teeth. After how many turns of each do they return together to their starting position?}{$\ppcm\pt{12;18}=36$ teeth: $3$ turns of the first, $2$ turns of the second.}{}
\end{fluencysheet}

\begin{fluencysheet}{Mon~9~Nov.~2026}{Drawing a line given by an equation}
\mxfq{Warm-up · Quick}{IMG}{For $y=-2x+3$, calculate $y$ when $x=0$ and then when $x=2$.}{$y=3$ then $y=-1$.}{}
\mxfq{Core · Method}{REG}{Give two points that allow you to draw the line $2x+y-4=0$.}{For $x=0$: $y=4$; for $x=2$: $y=0$. Points $\pt{0;4}$ and $\pt{2;0}$.}{}
\mxfq{Core + · Quick}{GEO}{What is the line with equation $x=3$? Give a direction vector. Does it have a gradient?}{The vertical line through $\pt{3;0}$; $\vc{0}{1}$ is a direction vector; it has no gradient; Cartesian equation $x-3=0$.}{Cannot be written $y=mx+p$.}
\mxfq{Review · Quick}{CMP}{Compare $\dfrac37$ and $\dfrac49$.}{Cross-multiplying: $3\times9=27<28=4\times7$, so $\dfrac37<\dfrac49$.}{}
\mxfq{Thread · Quick}{EQ}{Solve in $\R$: $5x-7=2x+8$.}{$x=5$}{Pitfall: collect the $x$ terms on one side, changing signs correctly.}
\end{fluencysheet}

\begin{fluencysheet}{Tue~10~Nov.~2026}{Finding the equation of a line: two points, a point and a direction vector, a point and a gradient}
\mxfq{Warm-up · Quick}{GEO}{Gradient of the line through $\pt{1;3}$ and $\pt{3;7}$?}{$\dfrac{7-3}{3-1}=2$.}{}
\mxfq{Core · Method}{GEO}{Slope-intercept equation of the line with gradient $2$ through $A\pt{1;3}$.}{$y-3=2(x-1)$, that is $y=2x+1$.}{Check: $2\times1+1=3$.}
\mxfq{Core + · Method}{GEO}{Line through $A\pt{1;2}$ with direction vector $\vec u\vc{2}{-1}$: Cartesian equation using the determinant of $\vect{AM}$ and $\vec u$.}{$M\pt{x;y}\in d\Leftrightarrow\deter{\vect{AM}}{\vec u}=0$: $(x-1)(-1)-(y-2)\times2=0$, that is $x+2y-5=0$.}{Check: $1+4-5=0$; $\vc{-2}{1}$ is indeed collinear to $\vc{2}{-1}$.}
\mxfq{Review · Quick}{PUI}{Simplify $2^5\times2^{-3}$.}{$2^{5-3}=2^2=4$.}{}
\mxfq{Thread · Table → formula}{REG}{An affine function $f$ satisfies: $f(0)=-1$, $f(1)=1$, $f(2)=3$. Give $f(x)$.}{Gradient $2$, y-intercept $-1$: $f(x)=2x-1$.}{}
\end{fluencysheet}

\begin{fluencysheet}{Thu~12~Nov.~2026}{Lines and affine functions: linking the representations}
\mxfq{Warm-up · Quick}{IMG}{$f(x)=-3x+4$. Calculate $f(2)$ and the preimage of $10$.}{$f(2)=-2$; $-3x+4=10\Leftrightarrow x=-2$.}{}
\mxfq{Core · Method}{REG}{Find the affine function $f$ such that $f(1)=1$ and $f(3)=5$.}{Gradient $\dfrac{5-1}{3-1}=2$; $f(x)=2x+p$ with $2+p=1$: $f(x)=2x-1$.}{}
\mxfq{Core + · True/false}{REG}{True or false? Justify. Every line in the plane is the graph of an affine function.}{False. Counter-example: the vertical line $x=3$ is not the graph of any function (the number $3$ would have infinitely many images).}{}
\mxfq{Review · Quick}{PRB}{A fair $6$-sided die is rolled. What is the probability of getting a multiple of $3$?}{Outcomes $3$ and $6$: $\dfrac26=\dfrac13$.}{}
\mxfq{Thread · Quick}{GEO}{$ABC$ is right-angled at $A$ with $AB=6$ cm and $AC=8$ cm. Calculate $BC$.}{$BC^2=36+64=100$, so $BC=10$ cm.}{}
\end{fluencysheet}

\begin{fluencysheet}{Mon~16~Nov.~2026}{Collinearity, parallelism, intersection of lines}
\mxfq{Warm-up · Quick}{SUB}{Does the point $\pt{3;5}$ lie on the line $y=2x-1$?}{$2\times3-1=5$: yes.}{}
\mxfq{Core · Method}{GEO}{Relative positions of $d:\ y=3x+1$ with $d':\ y=3x-2$, then with $d'':\ y=-x+1$.}{$d$ and $d'$: same gradient, different y-intercepts: strictly parallel; $d$ and $d''$: different gradients: intersecting.}{}
\mxfq{Core + · Method}{GEO}{Show in two ways that $A\pt{1;4}$, $B\pt{2;7}$ and $C\pt{4;13}$ are collinear.}{(i) All three satisfy $y=3x+1$. (ii) $\vect{AB}\vc{1}{3}$, $\vect{AC}\vc{3}{9}$: $\deter{\vect{AB}}{\vect{AC}}=9-9=0$.}{}
\mxfq{Review · Quick}{PRO}{On a map with scale $\dfrac{1}{\num{50000}}$, two towns are $4$ cm apart. What is the real distance?}{$4\times\num{50000}=\num{200000}$ cm $=2$ km.}{}
\mxfq{Thread · Method}{ALG}{Programme: choose a number $x$; add $3$; multiply by $2$; subtract $6$. Calculate the result for $x=5$, then write an expression in $x$ and simplify it. What do you notice?}{For $5$: $10$. Expression: $2(x+3)-6=2x$: the result is always twice the chosen number.}{Link A09: algebraic notation proves what numerical trials suggest.}
\end{fluencysheet}

\begin{fluencysheet}{Tue~17~Nov.~2026}{Intersection of two lines; Python collinearity function}
\mxfq{Warm-up · Quick}{EQ}{Solve $2x+1=-x+7$.}{$3x=6$, so $x=2$.}{}
\mxfq{Core · Method}{GEO}{Coordinates of the intersection of $y=2x+1$ and $y=-x+7$.}{$x=2$ then $y=2\times2+1=5$: $\pt{2;5}$.}{Check in the second equation: $-2+7=5$.}
\mxfq{Core + · Code reading}{ALG}{What do \code{alignes(0,0,1,2,2,4)} and \code{alignes(0,0,1,2,2,5)} return?\mxpseudo{def alignes(xa, ya, xb, yb, xc, yc):\\ \mxtab return (xb-xa)*(yc-ya) - (yb-ya)*(xc-xa) == 0}}{\code{True} then \code{False}: determinants $1\times4-2\times2=0$ and $1\times5-2\times2=1\neq0$.}{}
\mxfq{Review · Expand}{DEV}{Expand and simplify $3(x-2)$.}{$3x-6$.}{}
\mxfq{Thread · Prove}{LIT}{Show that the sum of three consecutive integers is a multiple of $3$.}{With $n-1,\ n,\ n+1$: sum $=3n$, a multiple of $3$.}{Choosing the middle term simplifies the calculation.}
\end{fluencysheet}
